
Give inductive proofs of the following 
\begin{enumerate}
\item $\forall x \in \Naturals, \; 3 \divides x^3-x$

\hint{
\begin{proof} (by induction)

{\bf Base case:} ($x=0$)

The base case is trivial since, when $x=0$, $x^3-x$ evaluates to $0$ and every natural number is a divisor of $0$.

{\bf Inductive step:} 

In the inductive step we want to show that $\forall k \in \Naturals, \; 3 \divides k^3 - k \; \implies \; 3 \divides (k+1)^3 - (k+1)$.  So, let $k$ be an arbitrary natural number such that  $3 \divides k^3 - k$.  

Consider the quantity $(k+1)^3 - (k+1)$,

\begin{gather*}
(k+1)^3 - (k+1) \\
= \; (k^3 + 3k^2+ 3k +1) - (k+1) \\
= \; (k^3 - k) + (3k^2 + 3k). 
\end{gather*}

By the inductive hypothesis (and the definition of divisibility), we know that there is an integer $q$ such that $k^3-k = 3q$.  So, by substitution, we obtain

\[ (k+1)^3 - (k+1) \; = \; 3q + 3k^2 + 3k \; = 3 \cdot (q + k^2 + k).\]

Finally, $q+k^2-k$ (being the sum of integers) is certainly an integer, so by the definition of 
divisibility, it follows that $(k+1)^3 - (k+1)$ is divisible by 3.

\end{proof}
}
\wbvfill

%\workbookpagebreak

\item $\forall x \in \Naturals, \; 3 \divides x^3+5x$

\hint{
	\begin{proof} (by induction)

{\bf Base case:} ($x=0$)

The base case is trivial since, when $x=0$, $x^3+5x$ evaluates to $0$ and every natural number is a divisor of $0$.

{\bf Inductive step:} 

In the inductive step we want to show that $\forall k \in \Naturals, \; 3 \divides k^3 + 5k \; \implies \; 3 \divides (k+1)^3 + 5(k+1)$.  So, let $k$ be an arbitrary natural number such that  $3 \divides k^3 + 5k$.  

Consider the quantity $(k+1)^3 + 5(k+1)$,

\begin{gather*}
(k+1)^3 + 5(k+1) \\
= \; (k^3 + 3k^2+ 3k +1) + 5(k+1) \\
= \; (k^3 + 5k) + (3k^2 + 3k + 6). 
\end{gather*}

By the inductive hypothesis (and the definition of divisibility), we know that there is an integer $q$ such that $k^3+5k = 3q$.  So, by substitution, we obtain

\[ (k+1)^3 + 5(k+1) \; = \; 3q + 3k^2 + 3k + 6\; = 3 \cdot (q + k^2 + k + 2).\]

Finally, $q+k^2+k+2$ (being the sum of integers) is certainly an integer, so by the definition of 
divisibility, it follows that $(k+1)^3 + 5(k+1)$ is divisible by 3.

\end{proof}
}
\wbvfill

\workbookpagebreak

\item $\forall x \in \Naturals, \; 11 \divides x^{11}+10x$

\hint{
	\begin{proof} (by induction)

{\bf Base case:} ($x=0$)

When $x=0$, $x^{11}+10x$ evaluates to $0$ and it is true that $11$ divides evenly into $0$.

{\bf Inductive step:} 

We want to show that $\forall k \in \Naturals, \; 11 \divides k^{11} + 10k \; \implies \; 11 \divides (k+1)^{11} + 10(k+1)$.  So, let $k$ be an arbitrary natural number such that  $11 \divides k^{11} + 10k$.  

Consider the quantity $(k+1)^{11} + 10(k+1)$.

\begin{gather*}
(k+1)^{11} + 10(k+1) \\
 = \; ( k^{11} + 11k^{10} + 55k^9 + 165k^8 + 330k^7 + 462k^6 + 462k^5 +330k^4 + 165 k^3 + 55k^2 + 11k + 1 ) + 10(k+1) \\
 = \; (k^11 + 10k) + ( 11k^{10} + 55k^9 + 165k^8 + 330k^7 + 462k^6 + 462k^5 +330k^4 + 165 k^3 + 55k^2 + 11k + 11 ) \\
 = \; (k^11 + 10k) + 11 \cdot (k^{10} + 5k^9 + 15k^8 + 30k^7 + 42k^6 + 42k^5 + 30k^4 + 15 k^3 + 5k^2 + k + 1 )
\end{gather*}

By the inductive hypothesis, we know that there is an integer $q$ such that $k^{11}+10k = 11q$.  So, by substitution, we obtain

\[ (k+1)^{11} + 10(k+1) \; = \;  11 \cdot ( q + k^{10} + 5k^9 + 15k^8 + 30k^7 + 42k^6 + 42k^5 + 30k^4 + 15 k^3 + 5k^2 + k + 1 ).\]

We've expressed $(k+1)^{11} + 10(k+1)$ as $11$ times a quanity that is clearly an integer (since it is the sum of powers and products of integers). Thus the inductive step is proven and the result follows by PMI.

\end{proof}

}
\wbvfill

%\workbookpagebreak

\item $\forall n \in \Naturals, \; 3 \divides 4^n-1$

\hint{
	\begin{proof} (by induction)

{\bf Base case:} ($n=0$)

Note that $4^0 - 1 = 0$ and it is certainly true that $3 \divides 0$.

{\bf Inductive step:} 

We want to show that $\forall k \in \Naturals, \; 3 \divides 4^k - 1 \; \implies \; 3 \divides 4^{k+1} - 1$.  So, let $k$ be an arbitrary natural number such that $3 \divides 4^k - 1$.  

Note that 

\begin{gather*}
4^{k+1} - 1 \\
= \; 4 \cdot 4^{k} - 1 \\
= \; (3+1) \cdot 4^{k} - 1 \\
= \; 3 \cdot 4^{k} + 4^{k} - 1 \\
\end{gather*}

By the inductive hypothesis, $4^{k} - 1 = 3q$ for some integer $q$ so

\begin{gather*}
4^{k+1} - 1 \\
= \; 3 \cdot 4^{k} + 3q \\
= \; 3 \cdot (4^{k} + q) \\
\end{gather*}

Finally, since $4^{k} + q$ is clearly an integer, the inductive conclusion is proved and the result follows by PMI.

\end{proof}
}



\wbvfill

\workbookpagebreak

\item $\forall n \in \Naturals, \; 6 \divides (3n^{2}+3n-12)$

\wbvfill

%\workbookpagebreak

\hint{
	\begin{proof} (by induction)

{\bf Base case:} ($n=0$)

Note that when $n=0$, the question becomes does $6$ divide $-12$, and the answer is ``yes.'' 


{\bf Inductive step:} 

We want to show that $\forall k \in \Naturals, \; 6 \divides (3k^{2}+3k-12)  \; \implies \; 6 \divides (3(k+1)^{2}+3(k+1)-12)$.  So, let $k$ be an arbitrary natural number such that $3 \divides (3k^{2}+3k-12)$.  

Note that 
\begin{gather*}
3(k+1)^{2}+3(k+1)-12
= \; 3(k^2 + 2k + 1) + 3(k+1) - 12 \\ 
= \; 3k^2 + 6k + 3 + 3k + 3 - 12 \\
= \; (3k^2 + 3k - 12) + (6k + 6)
\end{gather*}

Since $ 6 \divides (3k^2 + 3k - 12)$ (the inductive hypothesis) and $6$ clearly divides $(6k+6)$, it follow that 
$6 \divides 3(k+1)^{2}+3(k+1)-12$.  Thus the inductive step is proved and the result follows by mathematical induction.
}

\item $\forall n \in \Naturals, \; 5 \divides (n^{5}-5n^{3}+14n)$

\wbvfill

\workbookpagebreak

\hint{
	\begin{proof} (by induction)

{\bf Base case:} ($n=0$)

Note that when $n=0$, $n^5 - 5n^3 + 14n \; = \; 0$, and $5 \divides 0$.


{\bf Inductive step:} 

We want to show that $\forall k \in \Naturals, \; 5 \divides (k^5 - 5k^3 + 14k)  \; \implies \; 5 \divides ((k+1)^{5}-5(k+1)^3 + 14k)$.  So, let $k$ be an arbitrary natural number such that $5 \divides (k^5 - 5k^3 + 14k)$.  Consider the quantity $(k+1)^{5}-5(k+1)^3 + 14k$.

\begin{gather*}
(k+1)^{5} - 5(k+1)^3 + 14k \\
= \; k^5 + 5k^4 + 10k^3 + 10k^2 + 5k + 1 - 5(k^3 + 3k^2 + 3k + 1) + 14(k+1) \\
= \; (k^5 - 5k^3 +14k) + (5k^4 + 10k^3 + 10k^2 + 5k + 1 - 5(3k^2  + 3k + 1) + 14 \\
= \; (k^5 - 5k^3 +14k) + (5k^4 + 10k^3  - 5k^2 - 10k + 10) \\
=\;  (k^5 - 5k^3 +14k) + 5(k^4 + 2k^3 - k^2 - 2k + 2)
\end{gather*}

Since $k^4 + 2k^3 - k^2 - 2k + 2$ is clearly an integer, and the inductive hypothesis tells us that $k^5 - 5k^3 +14k$ is a multiple of $5$, it follows that 
$(k+1)^{5} - 5(k+1)^3 + 14k$  is a multiple of $5$ and so the inductive step is proved.  
\end{proof}
}

\item $\forall n \in \Naturals, \; 4 \divides (13^{n}+4n-1)$

\wbvfill

%\workbookpagebreak

\item $\forall n \in \Naturals, \; 7 \divides 8^n+6$

\wbvfill

\workbookpagebreak

\item $\forall n \in \Naturals, \; 6 \divides 2n^3 - 2n - 12$

\wbvfill

%\workbookpagebreak

\item $\forall n \geq 3 \in \Naturals, \; 3n^2+3n+1 < 2n^3$

\wbvfill

\workbookpagebreak

\item $\forall n > 3 \in \Naturals, \; n^3 < 3^n$

\wbvfill

%\workbookpagebreak

\item $\forall n \geq 3 \in \Naturals, \; n^{3}+3>n^{2}+3n+1$

\wbvfill

\workbookpagebreak

\item $\forall x \geq 4 \in \Naturals, \; x^22^x \leq 4^x$

\wbvfill

%\workbookpagebreak

\end{enumerate}


%% Emacs customization
%% 
%% Local Variables: ***
%% TeX-master: "GIAM-hw.tex" ***
%% comment-column:0 ***
%% comment-start: "%% "  ***
%% comment-end:"***" ***
%% End: ***

